\documentclass{article}
\usepackage[UTF8]{ctex}
\usepackage{amsmath,amssymb,amsfonts}
\usepackage{enumitem}
\usepackage{geometry}
\geometry{a4paper, margin=1in}
% === 错题本排版优化（不参与题目标识） ===
\setlength{\parindent}{0pt}
\setlength{\parskip}{0.35em plus 0.10em minus 0.08em}
\setlength{\emergencystretch}{3em}
\setlist{itemsep=0.25em, topsep=0.40em, parsep=0pt}
\allowdisplaybreaks[2]
% === 错题本排版优化结束 ===
\begin{document}
\section*{湘潭大学2024年高等代数真题}
\section*{187．湘潭大学2024年高等代数试卷}

\subsection{题目1}
1．计算 \(n\) 阶行列式

\[
D=\left|\begin{array}{cccccc}
x+y & x y & 0 & \cdots & 0 & 0 \\
1 & x+y & x y & \cdots & 0 & 0 \\
0 & 1 & x+y & \cdots & 0 & 0 \\
\vdots & \vdots & \vdots & & \vdots & \vdots \\
0 & 0 & 0 & \cdots & x+y & x y \\
0 & 0 & 0 & \cdots & 1 & x+y
\end{array}\right|
\]

\subsection{题目2}
2．设矩阵

\[
A=\left(\begin{array}{ccccc}
1 & 1 & 1 & \cdots & 1 \\
a_{1} & a_{2} & a_{3} & \cdots & a_{s} \\
a_{1}^{2} & a_{2}^{2} & a_{3}^{2} & \cdots & a_{s}^{2} \\
\vdots & \vdots & \vdots & & \vdots \\
a_{1}^{n-1} & a_{2}^{n-1} & a_{3}^{n-1} & \cdots & a_{s}^{n-1}
\end{array}\right)
\]

其中 \(a_{i} \neq a_{j}(i \neq j), \alpha_{1}, \alpha_{2}, \cdots, \alpha_{s}\) 分别为 \(A\) 的列向量，讨论 \(\alpha_{1}, \alpha_{2}, \cdots, \alpha_{s}\) 的线性相关性．

\subsection{题目3}
3．设方程组

\[
\left\{\begin{array}{l}
x_{1}+x_{2}+x_{3}+x_{4}=-1 \\
4 x_{1}+3 x_{2}+5 x_{3}-x_{4}=-1 \\
a x_{1}+x_{2}+3 x_{3}+b x_{4}=1
\end{array}\right.
\]

线性无关的解向量有 3 个．

\begin{itemize}
  \item[（1）] 求系数矩阵的秩．
  \item[（2）] 求 \(a, b\) 的值与方程组的通解。
\end{itemize}

\subsection{题目4}
4．设 \(A\) 为 \(n\) 阶方阵，证明：\(r\left(A^{n}\right)=r\left(A^{n+1}\right)\) 。

\subsection{题目5}
5．设 \(A\) 为 \(n\) 阶实对称矩阵，若 \(A-E\) 为正定矩阵。

\begin{itemize}
  \item[（1）] 证明：\(A\) 为正定矩阵。
  \item[（2）] 证明：\(E-A^{-1}\) 为正定矩阵。
\end{itemize}

\subsection{题目6}
6．设 \(f(x)=\left(x-a_{1}\right)\left(x-a_{2}\right) \cdots\left(x-a_{n}\right)-1\) ，其中 \(a_{1}, a_{2}, \cdots, a_{n}\) 为互不相同的整数．证明：\(f(x)\) 在有理数域上不可约．

\subsection{题目7}
7．设数域 \(P\) 上线性空间上的向量组 \(\alpha_{1}, \alpha_{2}, \cdots, \alpha_{n}\) 的秩为 \(r\) ，证明：使得 \(k_{1} \alpha_{1}+k_{2} \alpha_{2}+\cdots+k_{n} \alpha_{n}=0\)的 \(n\) 维向量 \(\left(k_{1}, k_{2}, \cdots, k_{n}\right)\) 全体构成的集合 \(V_{1}\) 是线性空间 \(P^{n}\) 的 \(n-r\) 维子空间．

\subsection{题目8}
8．设 \(\alpha_{1}, \alpha, \cdots, \alpha_{n}\) 为线性空间 \(V\) 的一组基，若 \(\left(\beta_{1}, \beta_{2}, \cdots, \beta_{m}\right)=\left(\alpha_{1}, \alpha_{2}, \cdots, \alpha_{n}\right) A\) ，其中 \(A\) 为 \(n \times m\)矩阵．证明： \(\operatorname{dim} L\left(\beta_{1}, \beta_{2}, \cdots, \beta_{m}\right)=\tau(A)\) 。

\subsection{题目9}
9．设 \(a, b, c\) 为实数，矩阵

\[
A=\left(\begin{array}{lll}
b & c & a \\
c & a & b \\
a & b & c
\end{array}\right), B=\left(\begin{array}{lll}
c & a & b \\
a & b & c \\
b & c & a
\end{array}\right), C=\left(\begin{array}{lll}
a & b & c \\
b & c & a \\
c & a & b
\end{array}\right) .
\]

\begin{itemize}
  \item[（1）] 证明：\(A, B, C\) 彼此相似．
  \item[（2）] 若 \(B C=C B\) ，则 \(A\) 至少有两个特征值为 0 ．
\end{itemize}

\subsection{题目10}
10．设 \(V_{1}\) 是有限维欧氏空间 \(V\) 的子空间，定义 \(V\) 到 \(V_{1}\) 的投影变换 \(\sigma\) 如下：\(\forall X \in V\) ，设 \(X=X_{1}+X_{2}\) ，\(X_{1} \in V_{1}, X_{2} \in V_{1}^{\perp}\) ，则 \(\sigma(X)=X_{1}\) 。

\begin{itemize}
  \item[（1）] 证明：\(\sigma\) 为 \(V\) 上的线性变换．
  \item[（2）] 证明：\(\sigma\) 是满足 \(\sigma=\sigma^{2}\) 的对称变换．
\end{itemize}
\end{document}
